\chapter{原型与初始基线}
本章训练与数据接口参考LeRobot项目\cite{lerobot}，Qwen模型加载参考官方实现\cite{qwen}。


基础平台先解决数据和动作接口是否可信，再讨论模型输出是否有意义。首轮Qwen适配、动作分块训练和跨模拟器流程都保留各自的条件。

\section{SO101基础平台与示教工具}

\label{stage:s01}

首先建立电脑与主从机械臂之间的连接，完成串口识别、校准、遥操作和示教记录。示教同时保存图像、关节状态与发送的动作，回放用于核对录制与执行接口。

\subsection{硬件、环境与可重复启动}\label{s01-detail-1}

最早的软件验证基于Ubuntu 22.04和RTX 4090工作站。厂家LeRobot 0.3.4工程保留，运行包装脚本隔离ROS2的Python搜索路径。环境名为hiwonder-so101；环境曾随旧目录迁移，部分入口仍引用旧绝对路径，因此采用明确解释器路径，不假定普通conda或pip入口都能工作。

\FloatBarrier\begingroup\small\setlength{\tabcolsep}{2.5pt}
\begin{longtable}[]{@{}
  >{\raggedright\arraybackslash}p{(\linewidth - 2\tabcolsep) * \real{0.3305}}
  >{\raggedright\arraybackslash}p{(\linewidth - 2\tabcolsep) * \real{0.6695}}@{}}
\caption{硬件、环境与可重复启动：设置与结果}\label{tab:s01-1}\tabularnewline
\toprule\noalign{}
\begin{minipage}[b]{\linewidth}\raggedright
项目
\end{minipage} & \begin{minipage}[b]{\linewidth}\raggedright
2026年10月12日记录
\end{minipage} \\
\midrule\noalign{}
\endfirsthead
\toprule\noalign{}
\begin{minipage}[b]{\linewidth}\raggedright
项目
\end{minipage} & \begin{minipage}[b]{\linewidth}\raggedright
2026年10月12日记录
\end{minipage} \\
\midrule\noalign{}
\endhead
\bottomrule\noalign{}
\endlastfoot
主臂 & 黑色SO101，作为遥操作输入 \\
从臂 & 白色SO101，执行目标关节位置 \\
黑臂稳定串口标识 & \path|usb-1a86_USB_Single_Serial_5CE5121669-if00| \\
白臂稳定串口标识 & \path|usb-1a86_USB_Single_Serial_5CE5122996-if00| \\
临时设备号 & 主臂ttyACM0、从臂ttyACM1；重插后可能变化 \\
摄像头 & icSpring，640\(\times\)480；当时图像节点为video0 \\
代码与校准 & SO101/code/lerobot；SO101/calibration/ \\
\end{longtable}\endgroup

串口角色由逐个连接识别，稳定标识用于启动；设备号不写成永久对应。主从臂电机校准值与保存文件做只读核对后，再由现场操作确认遥操作。这里建立的是硬件通信与动作接口基线。

\subsection{示教记录与回放接口}\label{s01-detail-2}

记录工具同时保存图像、实测关节状态、发送目标和主机时间戳。执行记录中的动作是实际下发的目标，不以主臂读数代替。相机先检查、记录准备后才开始采集；校准不一致则取消。一次完整示教已录制并回放，后续回放还加入起点对齐与分段确认。

不同时间有三种动作来源：人移动主臂产生目标、程序重放已录目标、模型依据当前图像与状态预测目标。三者共用串口与电机接口，但不共用``模型自主完成''的结论。固定轨迹的物体转移过程后来被反馈存在碰桌，因此不能保留为无碰撞通过案例。

\subsection{图像旋转与手腕偏置分别记录}\label{s01-detail-3}

画面旋转改变显示及模型所见像素，手腕偏置改变从臂关节目标，两者是独立变换。临时腕部映射为负90度时，原校准文件哈希保留；真实示教记录实际下发后的动作，因此训练和推理不能再次叠加同一个偏置。鼠标垫任务的原始图像以rotation=0记录，不能只依据网页显示方向推断训练图像方向。

\FloatBarrier\begingroup\small\setlength{\tabcolsep}{2.5pt}
\begin{longtable}[]{@{}
  >{\raggedright\arraybackslash}p{(\linewidth - 4\tabcolsep) * \real{0.2269}}
  >{\raggedright\arraybackslash}p{(\linewidth - 4\tabcolsep) * \real{0.3720}}
  >{\raggedright\arraybackslash}p{(\linewidth - 4\tabcolsep) * \real{0.4011}}@{}}
\caption{图像旋转与手腕偏置分别记录：设置与结果}\label{tab:s01-2}\tabularnewline
\toprule\noalign{}
\begin{minipage}[b]{\linewidth}\raggedright
信息
\end{minipage} & \begin{minipage}[b]{\linewidth}\raggedright
含义
\end{minipage} & \begin{minipage}[b]{\linewidth}\raggedright
需要保持一致的对象
\end{minipage} \\
\midrule\noalign{}
\endfirsthead
\toprule\noalign{}
\begin{minipage}[b]{\linewidth}\raggedright
信息
\end{minipage} & \begin{minipage}[b]{\linewidth}\raggedright
含义
\end{minipage} & \begin{minipage}[b]{\linewidth}\raggedright
需要保持一致的对象
\end{minipage} \\
\midrule\noalign{}
\endhead
\bottomrule\noalign{}
\endlastfoot
图像旋转 & 原始传感器像素到显示或模型像素的变换 & 图像、框、点、深度与相机内参 \\
腕部偏置 & 主从关节映射中的固定角度差 & 实测状态、发送动作和归一化统计 \\
电机校准 & 编码器零位与有效范围 & 硬件、保存文件及运行配置 \\
\end{longtable}\endgroup

编码器增量与角度的换算按4096计数对应一圈处理，即角度变化约为计数变化乘以360/4096。这个换算描述关节角增量，不等于末端空间高度或碰撞间隙。


\paragraph{阶段结论}完成过一条完整示教的录制及动作回放。后续腕部安装方向变化，又引入临时的负90度映射；这一映射必须与记录数据一同保存。 稳定设备标识、校准文件和角度约定需要随运行保存。轨迹回放说明记录与执行接口可用，不计为模型自主操作。

\stagecode{s01}

\FloatBarrier

\section{相机、场景理解与观察工作台}

\label{stage:s02}

把腕部相机、Qwen场景描述、YOLO目标检测和单目深度显示接入网页。模型完成一次推理后再取当时的新画面，操作人员可在等待期间阅读上一轮输出。

\subsection{连续观察的取帧方式}\label{s02-detail-1}

工作台由相机采集、检测与深度处理、场景描述三类工作组成。相机持续更新最新帧；Qwen完成一次推理后，才读取当时最新的图像开始下一轮。在新结果生成前，界面继续显示上一次结果，操作人员可以阅读。该设计没有把每帧都排队送入慢模型。

设图像在主机读取返回的时刻为 \(t_r\)，提交模型为 \(t_s\)，输出为 \(t_o\)，则取帧时的年龄为 \(t_s-t_r\)，输出对应画面的年龄为 \(t_o-t_r\)。这些时刻不包含传感器曝光时间；即使没有排队，较慢的生成仍会使输出对应数秒前的画面。

\begin{enumerate}
\def\labelenumi{\arabic{enumi}.}
\tightlist
\item
  相机线程写入最新RGB与序号。
\item
  快速视觉分支处理检测和深度，保留该阶段所用序号。
\item
  文字分支空闲后取最新帧，保存该帧和模型回复。
\item
  界面分别标明处理时间与画面年龄，避免把旧结果当成当前瞬间。
\end{enumerate}

\subsection{观察界面与运动权限}\label{s02-detail-2}

早期桌面窗口后来统一为网页，保留原RGB、深度或检测视图、文字场景观察。模型返回的结构化字段用于程序读取，界面可显示自然语言。观察结果不自动转成运动权限；在后续的有限技能系统中，只有通过明确任务范围检查的请求才会进入执行。

通用检测器没有覆盖特定耳机壳时，新增任务标注和专用检测模型。单目深度图可以显示相对远近，但未标定的米制输出不能直接用于限制机械臂下降。真实场景中曾出现观察到的空间距离与预测距离明显不符，因而深度分支保留为辅助观察，而不是唯一安全依据。

\subsection{共享功能与后续替换}\label{s02-detail-3}

网页也曾加入惯性测量单元的姿态显示，并通过相对初始参考姿态展示安装在连杆上的板卡。该功能属于观察基础设施，本文不展开另一条机器人控制课题的实验。之后的VLA网页把该区域换成动作模型输出与状态展示，但当串口被保持程序占用时，必须区分实时关节反馈与历史快照，不能把快照送入模型后称为完整实时闭环。


\stagefigure{figures/ch02/s02_web_workbench_preview.png}{早期网页工作台。图中保留当时界面，用于说明观察功能的组成。}{fig:s02-retained}

\paragraph{阶段结论}形成可以同时观察视频、检测结果和文字描述的原型。界面曾调整画面方向与深度色阶，后来收敛为网页工作台。 单目深度未经真实尺度标定，不能直接用作防碰撞距离。惯性传感器只作为共享观察功能介绍，其控制实验不并入本报告。

\stagecode{s02}

\FloatBarrier

\section{首轮Qwen适配与动作分块策略训练}

\label{stage:s03}

使用真实示教抽帧构造图像与指令用例，再训练Qwen的低秩适配器。动作侧采用ACT（动作分块变换器），在公开双臂交接数据上独立训练与仿真评测。

\subsection{Qwen 图像与指令评测}\label{s03-detail-1}

使用固定版本 Qwen3-VL-4B-Instruct，从既有 SO101 单条真实示教抽取 24 帧，分别提问``深蓝色塑胶壳''\,``蓝色耳机保护套''和不存在的``红色杯子''，共 72 条。输出目标可见性和 locate/observe 子任务，不输出关节命令。

参考标签来自单独一次自动图像审阅，未采用被测模型的答案作为标签，也未完成独立人工复核。24 帧同属一个视频，相关性高，结果只能称为暂定标签匹配率。

\FloatBarrier\begingroup\small\setlength{\tabcolsep}{2.5pt}
\begin{longtable}[]{@{}
  >{\raggedright\arraybackslash}p{(\linewidth - 2\tabcolsep) * \real{0.4967}}
  >{\raggedleft\arraybackslash}p{(\linewidth - 2\tabcolsep) * \real{0.5033}}@{}}
\caption{Qwen 图像与指令评测：设置与结果}\label{tab:s03-1}\tabularnewline
\toprule\noalign{}
\begin{minipage}[b]{\linewidth}\raggedright
指标
\end{minipage} & \begin{minipage}[b]{\linewidth}\raggedleft
结果
\end{minipage} \\
\midrule\noalign{}
\endfirsthead
\toprule\noalign{}
\begin{minipage}[b]{\linewidth}\raggedright
指标
\end{minipage} & \begin{minipage}[b]{\linewidth}\raggedleft
结果
\end{minipage} \\
\midrule\noalign{}
\endhead
\bottomrule\noalign{}
\endlastfoot
JSON 格式有效 & 72/72 \\
目标可见性及子任务匹配 & 66/72（91.67\%） \\
深蓝色塑胶壳 & 23/24；正例召回 12/13 \\
蓝色耳机保护套 & 19/24；正例召回 8/13 \\
不存在的红色杯子 & 24/24 正确拒绝 \\
耗时中位数 / P95 & 0.380 / 0.566 秒 \\
\end{longtable}\endgroup

恒定回答``不可见''能匹配 46/72，说明类别分布偏斜，应同时看正例召回。换一种描述后漏识别增加，不能以单一总指标断言稳定理解。耗时针对短结构化回答，不直接对比工作台较长的自然语言描述。

\stagefigure{figures/ch02/s03_qwen_instruction_08ff905c.png}{Qwen 评测}{fig:s03-detail-1}

\subsection{Qwen LoRA 首轮训练}\label{s03-detail-2}

按图像帧划分 48/12/12 条训练/验证/测试用例；同一帧的不同问法在同一分组。仍来自同一视频，不是跨场景验证。冻结视觉模块，仅对语言 self-attention 的 q\_proj/v\_proj 加 LoRA：r=8、alpha=16、dropout=0.05、学习率 5e-5、120 次更新、梯度累积 2，训练参数 2,949,120。仅监督回答 token，保留输入前缀核对。

\FloatBarrier\begingroup\small\setlength{\tabcolsep}{2.5pt}
\begin{longtable}[]{@{}
  >{\raggedright\arraybackslash}p{(\linewidth - 4\tabcolsep) * \real{0.3370}}
  >{\raggedleft\arraybackslash}p{(\linewidth - 4\tabcolsep) * \real{0.3315}}
  >{\raggedleft\arraybackslash}p{(\linewidth - 4\tabcolsep) * \real{0.3315}}@{}}
\caption{Qwen LoRA 首轮训练：设置与结果}\label{tab:s03-2}\tabularnewline
\toprule\noalign{}
\begin{minipage}[b]{\linewidth}\raggedright
分组
\end{minipage} & \begin{minipage}[b]{\linewidth}\raggedleft
基座
\end{minipage} & \begin{minipage}[b]{\linewidth}\raggedleft
LoRA
\end{minipage} \\
\midrule\noalign{}
\endfirsthead
\toprule\noalign{}
\begin{minipage}[b]{\linewidth}\raggedright
分组
\end{minipage} & \begin{minipage}[b]{\linewidth}\raggedleft
基座
\end{minipage} & \begin{minipage}[b]{\linewidth}\raggedleft
LoRA
\end{minipage} \\
\midrule\noalign{}
\endhead
\bottomrule\noalign{}
\endlastfoot
验证 12 条 & 11/12 & 12/12 \\
测试 12 条 & 11/12 & 11/12 \\
\end{longtable}\endgroup

测试未改善，训练损失极低可能反映记忆，不能宣称泛化提升。另一次缩写提示词重新加载测试出现字段遗漏及子任务不一致；使用训练协议的完整提示词重新加载通过。两份结果均保留，说明适配器存在提示词敏感性。

\stagefigure{figures/ch02/s03_qwen_lora_comparison_3da6789e.png}{LoRA 对比}{fig:s03-detail-2}

\subsection{ACT 训练与实际物理仿真评测}\label{s03-detail-3}

采用公开 ALOHA 双臂方块交接示教：50 个 episode、20,000 帧、50Hz、480\(\times\)640 RGB、14 维双臂关节/夹爪状态与动作。逐 episode 核对图像与动作帧数，按 episode 40/5/5 拆分；自身训练的归一化统计仅来自训练集，未来动作序列不跨 episode。

自训练 ACT 从随机策略参数开始，仅 ResNet18 使用 ImageNet 初始化，\textbf{没有加载公开 ACT 策略权重}。chunk=100、batch=8、学习率 1e-4、backbone 1e-5、seed=2026，共 5,000 更新，记录训练耗时 236.3 秒。验证与测试各取 64 个分层样本评估归一化动作 L1。

\begin{itemize}
\tightlist
\item
  验证 L1：0.9928 \(\rightarrow\) 0.2844；最佳验证 0.2714。
\item
  最终 模型检查点 的测试 L1：0.2839。
\item
  离线误差下降不等于抓取成功，另做实际 rollout。
\end{itemize}

\stagefigure{figures/ch02/s03_act_learning_604e4f8c.png}{ACT 曲线}{fig:s03-detail-3}

在 \path|gym_aloha/AlohaTransferCube-v0| 实际运行策略，固定 seeds 1000--1009，每次最多 400 步；最大任务 reward\(\geq\)4 作为方块交接成功。每次重置策略动作缓存，模型输入当前 RGB 和关节状态，保存逐次动作、奖励、视频。

\FloatBarrier\begingroup\small\setlength{\tabcolsep}{2.5pt}
\begin{longtable}[]{@{}
  >{\raggedright\arraybackslash}p{(\linewidth - 4\tabcolsep) * \real{0.2875}}
  >{\raggedleft\arraybackslash}p{(\linewidth - 4\tabcolsep) * \real{0.2325}}
  >{\raggedright\arraybackslash}p{(\linewidth - 4\tabcolsep) * \real{0.4800}}@{}}
\caption{ACT 训练与实际物理仿真评测：设置与结果}\label{tab:s03-3}\tabularnewline
\toprule\noalign{}
\begin{minipage}[b]{\linewidth}\raggedright
策略来源
\end{minipage} & \begin{minipage}[b]{\linewidth}\raggedleft
成功
\end{minipage} & \begin{minipage}[b]{\linewidth}\raggedright
说明
\end{minipage} \\
\midrule\noalign{}
\endfirsthead
\toprule\noalign{}
\begin{minipage}[b]{\linewidth}\raggedright
策略来源
\end{minipage} & \begin{minipage}[b]{\linewidth}\raggedleft
成功
\end{minipage} & \begin{minipage}[b]{\linewidth}\raggedright
说明
\end{minipage} \\
\midrule\noalign{}
\endhead
\bottomrule\noalign{}
\endlastfoot
公开预训练 ACT & 7/10（70\%） & 外部对照，训练使用了全部 50 episode，不能算自主训练成果 \\
自训练 ACT 最终权重 & 2/10（20\%） & 首轮 baseline，尚不稳定；两个失败试次达到 reward 2，但未完成交接 \\
\end{longtable}\endgroup

同样种子用于配对比较；样本只有 10 次，未给置信稳定性或真机迁移结论。公开策略与自训练数据量/训练配置不完全一致，此处是参考对照而非公平消融。

\stagefigure{figures/ch02/s03_act_rollout_comparison.png}{成功率对比}{fig:s03-detail-4}

\stagefigure{figures/ch02/s03_rollout_contact_sheet_f2db237e.jpg}{实际仿真关键帧}{fig:s03-detail-5}

另存单试次时延检查，不混入 10 次成功率统计。\path|select_action| 多数调用只取缓存动作，不能把其亚毫秒中位数当成整块网络推理耗时；完整 chunk 单独记录于 \path|*_act_latency_001/latencies.csv|，样本很少，含首次冷启动，不作为稳定性能基准。

\subsection{有限分层链路与失败修复}\label{s03-detail-4}

该阶段集成明确使用\textbf{基座 Qwen + 公开 ACT}，未把首轮 LoRA 或低成功率的自训练 ACT 当成集成模型。Qwen 看仿真初始 RGB 和指令，选择唯一已实现技能 transfer\_red\_cube 或 observe；ACT 执行后用仿真奖励记录结果。这里是单次高层选择和底层视觉反馈执行，尚未实现 Qwen 连续监控、自动重试或多步规划。

测试 5 条红方块交接指令、2 条蓝球指令，seeds1010--1016。初版模型技能选择 5/7 正确；两条蓝球指令均误选红方块技能并执行，说明仅靠提示词拒绝并不可靠。

增加独立于模型输出和测试标签的有限指令语法授权：只接受明确红方块抓取指令，否定或未支持目标进入 observe。相同用例复测：模型自身仍 5/7 正确，但有效派发 7/7 正确，蓝球误执行从 2/2 降至 0/2；红方块执行成功仍 3/5。\textbf{7/7 是受约束派发结果，不是 Qwen 准确率提升。} 指令语法覆盖范围有限，不能代替通用语言理解。


\paragraph{阶段结论}首轮Qwen测试为72个图像指令组合；ACT建立了自训练、公开权重和实际仿真回合之间的对照。 同一视频抽取的画面相关性较强。早期匹配率和动作回合数应按原报告分别列出，不合并成通用机器人成功率。

\stagecode{s03}

\FloatBarrier

\section{训练扩展、深度定位与初步流程衔接}

\label{stage:s04}

继续训练ACT，增加模拟深度、场景划分和Qwen定位适配，并测试语言部分INT8（8位整数）与键值缓存。导航和机械臂规划也开始接入任务流程。

\subsection{ACT追加训练与更可靠的评测}\label{s04-detail-1}

延续自训练5,000更新的权重，再训练15,000更新，累计20,000。使用相同公开示教、40/5/5 episode划分、训练集统计、chunk100、batch8、学习率1e-4/backbone1e-5。\textbf{优化器重新初始化，不是包含优化器状态的无缝续训；没有载入公开ACT权重。} 该阶段训练耗时798.2秒，期间部分其他GPU实验同时运行，不能作为纯训练吞吐基准。

追加训练的最佳验证L1为0.2588，最终权重验证0.2850、测试64抽样帧L1为0.2909。最终权重的离线误差没有改善，因此以\textbf{验证集}选择best checkpoint，不能把最后一次更新自动当成最好模型。

\stagefigure{figures/ch02/s04_act_continuation_53e2c29e.png}{继续训练的验证曲线}{fig:s04-detail-1}

使用另5个调参种子2000--2004，比较执行10/25/100动作后重新观察：成功分别2/5、1/5、4/5，选100。选择规则预设为成功数、平均最高奖励、相同则较长动作间隔。新测试种子3000--3029不参与选择；这里没有证明更频繁更新必然更好。

\FloatBarrier\begingroup\small\setlength{\tabcolsep}{2.5pt}
\begin{longtable}[]{@{}
  >{\raggedright\arraybackslash}p{(\linewidth - 4\tabcolsep) * \real{0.3103}}
  >{\raggedleft\arraybackslash}p{(\linewidth - 4\tabcolsep) * \real{0.3130}}
  >{\raggedright\arraybackslash}p{(\linewidth - 4\tabcolsep) * \real{0.3767}}@{}}
\caption{ACT追加训练与更可靠的评测：设置与结果}\label{tab:s04-1}\tabularnewline
\toprule\noalign{}
\begin{minipage}[b]{\linewidth}\raggedright
策略
\end{minipage} & \begin{minipage}[b]{\linewidth}\raggedleft
同30个测试种子成功
\end{minipage} & \begin{minipage}[b]{\linewidth}\raggedright
说明
\end{minipage} \\
\midrule\noalign{}
\endfirsthead
\toprule\noalign{}
\begin{minipage}[b]{\linewidth}\raggedright
策略
\end{minipage} & \begin{minipage}[b]{\linewidth}\raggedleft
同30个测试种子成功
\end{minipage} & \begin{minipage}[b]{\linewidth}\raggedright
说明
\end{minipage} \\
\midrule\noalign{}
\endhead
\bottomrule\noalign{}
\endlastfoot
原自训练5k最终权重 & 8/30（26.67\%） & 与上一份交付保持一致的基线 \\
原自训练5k最佳验证权重 & 8/30（26.67\%） & 补充相同验证选模规则的对照 \\
追加至20k最佳验证权重 & 16/30（53.33\%） & 同样动作间隔100；提高8次，尚不稳定 \\
公开预训练ACT & 25/30（83.33\%） & 外部参考，不是自主训练成果 \\
\end{longtable}\endgroup

\stagefigure{figures/ch02/s04_act_test30.png}{30次实际rollout结果}{fig:s04-detail-2}

相对5k最佳验证模型，11个种子由失败变成功，3个由成功变失败；并非所有条件都改善。图中为95\% Wilson区间，30次仍有较大不确定性。公开模型训练量及配置不同，不能作为公平训练消融。原10次结果与新30次用不同种子，分别报告，不用30次结果改写旧数据。

每次实际MuJoCo运行400步，保存动作、奖励与视频，以最高任务reward\(\geq\)4定义交接成功。离线验证误差不能代替这些成功率。多种子只覆盖当前模拟器的方块初始位置分布，不代表真实机器人或其他任务泛化。

\subsection{独立于单条真实视频的RGB与深度评测}\label{s04-detail-2}

新生成30个MuJoCo场景，seeds4000--4029；红/蓝/绿方块各10场景，改变方块尺寸、位置与相机，生成120条目标指令及否定指令。每场景仍仅一个可见方块，背景结构基本一致，不是复杂真实场景。

标签来源为同一几何场景中目标材质黑/白两次渲染的差分可见像素，得到真实目标框；深度来自模拟器光学深度渲染，按米保存。标签不由被测Qwen生成。首次调用dm\_control分割渲染出现索引错误，改用成对材质差分，明确保留方案与生成代码。

\stagefigure{figures/ch02/s04_contact_sheet_b6636723.jpg}{数据示例}{fig:s04-detail-3}

Qwen输入RGB和指令，输出目标可见性、归一化bbox与技能建议；后端根据预测框中心5\(\times\)5深度块取中位数。深度进入后端目标描述，\textbf{并未作为额外网络输入，也未实现相机到机械臂/地图外参}。

\FloatBarrier\begingroup\small\setlength{\tabcolsep}{2.5pt}
\begin{longtable}[]{@{}
  >{\raggedright\arraybackslash}p{(\linewidth - 2\tabcolsep) * \real{0.6047}}
  >{\raggedleft\arraybackslash}p{(\linewidth - 2\tabcolsep) * \real{0.3953}}@{}}
\caption{独立于单条真实视频的RGB与深度评测：设置与结果}\label{tab:s04-2}\tabularnewline
\toprule\noalign{}
\begin{minipage}[b]{\linewidth}\raggedright
基座Qwen指标
\end{minipage} & \begin{minipage}[b]{\linewidth}\raggedleft
120用例结果
\end{minipage} \\
\midrule\noalign{}
\endfirsthead
\toprule\noalign{}
\begin{minipage}[b]{\linewidth}\raggedright
基座Qwen指标
\end{minipage} & \begin{minipage}[b]{\linewidth}\raggedleft
120用例结果
\end{minipage} \\
\midrule\noalign{}
\endhead
\bottomrule\noalign{}
\endlastfoot
格式、目标可见性、技能建议匹配 & 各120/120 \\
正例定位IoU\(\geq\)0.5 & 58/60 \\
已检测正例平均IoU & 0.7032 \\
框中心深度平均绝对误差 & 0.00057m \\
否定指令技能匹配 & 30/30 \\
\end{longtable}\endgroup

深度误差很小，主要因模拟深度理想、方块可见、中心点通常落在目标面，不可当成真实深度精度/标定结果。120条来自30张图，存在共享图像与场景结构相关性，不能当成120个独立真实场景。技能匹配是接口建议测试；真正执行还须通过更严格的完整任务授权，不能仅因``红色可抓取''便执行完整交接。

\subsection{新LoRA：按场景分组的定位适配}\label{s04-detail-3}

18/6/6场景划分训练/验证/测试，对应72/24/24指令。同图所有问法在同一组，不与训练混用。标签来自上述模拟器差分框及明确技能规则，不采用模型自标注。

冻结视觉模型，仅语言q\_proj/v\_proj使用r8、alpha16、dropout0.05、学习率5e-5，训练240更新，2,949,120可训练参数。仅监督回答token，核对输入前缀一致。独立进程重新加载保存适配器后，评估未参与训练的48条验证/测试用例。

\FloatBarrier\begingroup\small\setlength{\tabcolsep}{2.5pt}
\begin{longtable}[]{@{}
  >{\raggedright\arraybackslash}p{(\linewidth - 6\tabcolsep) * \real{0.2485}}
  >{\raggedleft\arraybackslash}p{(\linewidth - 6\tabcolsep) * \real{0.2308}}
  >{\raggedleft\arraybackslash}p{(\linewidth - 6\tabcolsep) * \real{0.2327}}
  >{\raggedright\arraybackslash}p{(\linewidth - 6\tabcolsep) * \real{0.2879}}@{}}
\caption{新LoRA：按场景分组的定位适配：设置与结果}\label{tab:s04-3}\tabularnewline
\toprule\noalign{}
\begin{minipage}[b]{\linewidth}\raggedright
分组
\end{minipage} & \begin{minipage}[b]{\linewidth}\raggedleft
基座平均IoU
\end{minipage} & \begin{minipage}[b]{\linewidth}\raggedleft
LoRA平均IoU
\end{minipage} & \begin{minipage}[b]{\linewidth}\raggedright
可见性/技能匹配
\end{minipage} \\
\midrule\noalign{}
\endfirsthead
\toprule\noalign{}
\begin{minipage}[b]{\linewidth}\raggedright
分组
\end{minipage} & \begin{minipage}[b]{\linewidth}\raggedleft
基座平均IoU
\end{minipage} & \begin{minipage}[b]{\linewidth}\raggedleft
LoRA平均IoU
\end{minipage} & \begin{minipage}[b]{\linewidth}\raggedright
可见性/技能匹配
\end{minipage} \\
\midrule\noalign{}
\endhead
\bottomrule\noalign{}
\endlastfoot
验证12个正例 & 0.7281 & 0.8644 & 双方均24/24 \\
测试12个正例 & 0.7254 & 0.8275 & 双方均24/24 \\
\end{longtable}\endgroup

\stagefigure{figures/ch02/s04_lora_localization_5a072d09.png}{定位适配对照}{fig:s04-detail-4}

这轮的定位框在模拟场景测试上改善，语义准确率本来已满分，未表现出额外提升。样本小、单方块、同模拟器，训练损失极低可能含记忆；不宣称真实世界或跨模拟器泛化。与上一轮单视频LoRA是两套独立数据/接口，旧结果继续保留，不能用新结果替代旧测试。

\subsection{本机INT8与KV Cache实际对照}\label{s04-detail-4}

固定40条用例、相同模型/提示词/生成限制，训练及该阶段其他GPU模型结束后依次运行三组；每组独立进程加载，耗时分位数排除首次生成。未作多轮交错重复，因此数值是一次受控对照，不是稳定生产时延。

\FloatBarrier\begingroup\small\setlength{\tabcolsep}{2.5pt}
\begin{longtable}[]{@{}
  >{\raggedright\arraybackslash}p{(\linewidth - 8\tabcolsep) * \real{0.2452}}
  >{\raggedleft\arraybackslash}p{(\linewidth - 8\tabcolsep) * \real{0.1968}}
  >{\raggedleft\arraybackslash}p{(\linewidth - 8\tabcolsep) * \real{0.1860}}
  >{\raggedleft\arraybackslash}p{(\linewidth - 8\tabcolsep) * \real{0.1860}}
  >{\raggedright\arraybackslash}p{(\linewidth - 8\tabcolsep) * \real{0.1860}}@{}}
\caption{本机INT8与KV Cache实际对照：设置与结果}\label{tab:s04-4}\tabularnewline
\toprule\noalign{}
\begin{minipage}[b]{\linewidth}\raggedright
配置
\end{minipage} & \begin{minipage}[b]{\linewidth}\raggedleft
加载后已分配显存GiB
\end{minipage} & \begin{minipage}[b]{\linewidth}\raggedleft
单次推理中位秒
\end{minipage} & \begin{minipage}[b]{\linewidth}\raggedleft
P95秒
\end{minipage} & \begin{minipage}[b]{\linewidth}\raggedright
目标/技能匹配
\end{minipage} \\
\midrule\noalign{}
\endfirsthead
\toprule\noalign{}
\begin{minipage}[b]{\linewidth}\raggedright
配置
\end{minipage} & \begin{minipage}[b]{\linewidth}\raggedleft
加载后已分配显存GiB
\end{minipage} & \begin{minipage}[b]{\linewidth}\raggedleft
单次推理中位秒
\end{minipage} & \begin{minipage}[b]{\linewidth}\raggedleft
P95秒
\end{minipage} & \begin{minipage}[b]{\linewidth}\raggedright
目标/技能匹配
\end{minipage} \\
\midrule\noalign{}
\endhead
\bottomrule\noalign{}
\endlastfoot
BF16＋KV Cache & 8.266 & 0.456 & 0.713 & 40/40 \\
语言INT8＋KV Cache & 4.891 & 1.920 & 3.082 & 40/40 \\
BF16关闭KV Cache & 8.266 & 1.432 & 2.319 & 40/40 \\
\end{longtable}\endgroup

\stagefigure{figures/ch02/s04_quantization_cache_a85968c4.png}{量化及缓存}{fig:s04-detail-5}

INT8转换252个语言Linear，视觉和lm\_head未量化；显存下降约40.8\%，但本机慢约4.2倍。关闭KV Cache也明显更慢。因此当前工作台仍保留BF16＋缓存，不把省显存当成加速。量化推理会发生BF16\(\rightarrow\)FP16的激活转换，记录在日志里。\href{https://huggingface.co/docs/transformers/v4.57.0/quantization/bitsandbytes}{Transformers量化说明}

实际混合精度量化权重约5GB，已本地保存，新的进程从该目录直接加载并通过单例检查。它是\textbf{4090上验证的PC量化模型}，不是已验证的Jetson模型；不能推断Orin速度、功耗或ARM算子兼容性。

\subsection{Nav2、MoveIt与Qwen任务流程}\label{s04-detail-5}

Nav2使用实际Gazebo11/TurtleBot3、既有地图、AMCL、规划/控制节点；进入桌面站与返回两个action均完成。默认位置容差约0.25m，action成功不等于毫米级站点定位，odom记录不能直接当成map坐标误差。

MoveIt2使用独立domain63、Panda mock\_components，3个合法关节目标均规划并模拟执行；另一个超界目标未找到有效目标状态，返回错误-27，没有执行。保留实际规划点和mock关节反馈。\textbf{这是模拟控制接口，未包含接触物理抓取}；MoveGroup返回executed\_trajectory为空，不能捏造实际执行轨迹，反馈记录单独记录。

进一步实际跑通：

\begin{enumerate}
\def\labelenumi{\arabic{enumi}.}
\tightlist
\item
  当前MuJoCo初始RGB
\item
  基座Qwen建议
\item
  完整任务语法与时效授权
\item
  Nav2到固定桌面站
\item
  公开ACT方块交接
\item
  读取实际仿真奖励
\item
  Nav2返回
\end{enumerate}

一次流程通过，原始图像、Qwen回复、授权、Nav2 action、ACT视频/奖励、返回和真实仿真rosbag均保存。使用\textbf{公开ACT及已知成功seed1003的接口冒烟测试}，不是自训练模型成功率，也不是未见种子泛化测试。

\textbf{Gazebo底盘和MuJoCo双臂位于独立模拟器，没有共同机器人模型/物理世界/外参；MoveIt是另一个独立规划接口。} 当前验证任务流程与模块交接，不称``一台移动机械臂完成同场景导航抓取''。Qwen只做一次高层建议，尚无持续完成状态视觉判断、重试或自主恢复。

\stagefigure{figures/ch02/s04_navigation_trace_2f625cf4.png}{实际导航里程计}{fig:s04-detail-6}

\subsection{可代测接口、失败记录及研究草稿}\label{s04-detail-6}

完整流程录制仿真ROS2 bag，共2364条消息，包含odom、scan、cmd\_vel、clock、tf与tf\_static；SQLite完整性检查通过。可作为后续代测收集/解析接口的演练，来源仍为模拟器。

任务授权只接受完整支持的交接指令，不匹配目标替换、否定、附加任务、只抓取请求、过期结果或取消；任一阶段失败停止后续派发。15项软件检查通过，不能当作15次实际机器人鲁棒性实验。

首次Nav2在生命周期激活前发goal被拒绝；第二次在AMCL初始定位之前等待导航active形成启动依赖，最终超时。已改为AMCL active\(\rightarrow\)initialpose\(\rightarrow\)bt\_navigator active\(\rightarrow\)action，两个失败目录保留。shutdown日志中进程SIGINT/-2不是目标阶段的成功/失败判据。

技术论文草稿已围绕该阶段真实数据写成，注明工程复现、样本和创新边界。

下一步应优先同场景机器人与坐标系衔接、复杂目标及独立真实数据，再开展Jetson和企业代测；该阶段没有将可离线完成的准备工作写成未实际完成的硬件成果。


\paragraph{阶段结论}形成了Nav2导航、MoveIt规划与ACT执行的初步衔接，以及按验证集选择检查点的训练记录。 这一时期导航与抓取处于不同模拟环境，MoveIt执行使用模拟接口；它还不是同一物理世界里的移动操作。

\stagecode{s04}

\FloatBarrier
